\section{Tropical feasibility and further questions}\label{sec:closing}

The prices in the LP algorithm enter a comparison by addition to a logarithmic
magnitude. A system of tropical linear inequalities makes such comparisons
directly. Akian, Gaubert and Guterman~\cite{AGG} related its feasibility
to the values of a mean-payoff game (see also~\cite{MSS,DG}).
Theorem~\ref{thm:games} thus gives one further algorithmic consequence.

\begin{corollary}[Tropical linear feasibility]\label{cor:tropical}
Let \(a_{ri},b_{ri}\) be explicitly given entries in
\(\RR\cup\{-\infty\}\) of two \(m\times n\) matrices.
For the system
\[
                \max_i(a_{ri}+x_i)\ \le\ \max_j(b_{rj}+x_j)
                  \qquad\text{for every row \(r\)},
\]
a deterministic strongly polynomial algorithm finds a solution in
\((\RR\cup\{-\infty\})^n\) finite at every coordinate that can be
finite in any solution. An empty maximum equals \(-\infty\).
\end{corollary}

\begin{proof}
Use the game of~\cite{AGG}. At coordinate \(i\), Min may choose any row
\(r\) with \(a_{ri}\) finite. At row \(r\), Max may choose any coordinate
\(j\) with \(b_{rj}\) finite. The two successive edges have rewards
\(-a_{ri}\) and \(b_{rj}\). Give a vertex with no move a self-loop of
reward zero if it belongs to Min, and of reward \(-1\) if it belongs to Max.

A coordinate \(i\) finite in a solution has nonnegative game value.
After Min chooses a row there, Max takes a term attaining the right-hand
maximum. Its coordinate \(j\) is finite and
\(-a_{ri}+b_{rj}\ge x_i-x_j\). Summing over the play telescopes.
The finite coordinates form a finite set, so the bounded endpoint
contribution vanishes in the mean payoff.

Conversely, compute the values and fix an optimal policy for Max by
Theorem~\ref{thm:games}. From a coordinate of nonnegative value, Min can
reach no negative-weight cycle under this policy, since Min could then
repeat the cycle. In particular, every reachable Max vertex has a move
in the original game. At each coordinate reachable from a nonnegative-value
start, set \(x_i\) to the least total weight of a walk starting at \(i\)
with Max's policy fixed, allowing the empty walk. Erasing cycles shows
that a simple walk attains the minimum. At all other coordinates put
\(x_i=-\infty\). For each permitted pair of moves \(i,r,j\) from a
finite coordinate, concatenating the two moves with a minimizing walk
from \(j\) gives \(x_i\le -a_{ri}+b_{rj}+x_j\). Hence the row
inequalities hold. Dynamic programming over walks of length at most
the number of vertices computes the minima in strongly polynomial time.
\end{proof}

A constant term in a maximum can be homogenized by adding one coordinate.
The original system is feasible precisely when this coordinate can be finite.
In that case subtract its value from every coordinate to recover a solution.

\paragraph{Further questions.}
Allamigeon, Benchimol, Gaubert and Joswig~\cite{ABGJgame} showed how a
strongly polynomial combinatorial simplex algorithm could yield strongly
polynomial time for mean-payoff games. The algorithm of Theorem~\ref{thm:lp}
uses quantitative comparisons of minors as well as their signs. Some
comparisons in Section~\ref{sec:compression} have degree exponential in
the dimensions. Can a strongly polynomial LP algorithm use predicates of
polynomial degree throughout?

Theorem~\ref{thm:lp} takes the constraints as explicit input. For a
polyhedron given by a separation oracle, known exact algorithms also
involve geometric parameters~\cite{DVZ26}. A further question concerns
transition probabilities in games. Theorem~\ref{thm:games} permits
arbitrary real probabilities for one player. For two players its step
bound still depends on the bit lengths of the probabilities.
Can that dependence be removed?
